% Original aggregate-model chapters, retained outside the new main text.
\section{Classical Model}\label{classical-model}


\begin{greenblock}
Learning points for 23-2 Course: the goal is to understand:

\begin{itemize}
\tightlist
  \item motivations for contruction of a model
  \item the common approach to analyze a model
  \item how models evolve over time
  \item target function and constraint
\end{itemize}

\end{greenblock}

In a perfectly competitive market, the economic system can be described as\footnote{See in Sargint, 1987}:

\begin{align}
&\text{AS:}
\begin{cases}
Y=A\cdot F(K,N)\text{ (Only N is endogenously chosen; Inada conditions)}
\\
\max \pi \Rightarrow \frac{W}{P}=F_{N}(K,N) \\
N=N(\frac{W}{P}),N'>0\text{ (ad hoc, assuming labor market clearing)}\end{cases}
\\
&\text{AD: }
\begin{cases}
Y=C+I+G+\delta K \\
C=C(Y,T-\pi) \\
I=I(q-1), I'>0, q\equiv \frac{F_{k}-(r+\delta-\pi)}{r-\pi} \\
\frac{M}{P}=m(Y,r),m_{1}>0,m_{2}<0
\end{cases}
\end{align}

\begin{blueblock}
\begin{note}[Exogeneous variables in this model would be]
\end{note}
\begin{itemize}
\tightlist
  \item In AS: $K$
  \item In AD: $K, G, M, \pi, \delta, T$
\end{itemize}
\end{blueblock}


At the equilibrium, the endogenous variables could be determined jointly by multiple exo variables. e.g.

$$
\text{AS}
 \Rightarrow
\begin{cases}
Y^{*}=f(k) \\
N^{*}=g(k) \\
(\frac{W}{P})^{*}=h(k)
\end{cases}
$$


\begin{redblock}
\begin{question}[how to calculate $N$ at the equilibrium?]
\end{question}
By using total derivatives, we could have:
$$\begin{cases}
&dY=dY(dK)
\\&d \frac{W}{P}=d \frac{W}{P}(dK)
\\&dW=dW(dK)
\\&dP=dP(dK)
\end{cases}$$
\end{redblock}

\begin{blueblock}
\begin{note}
All a model can contain are:
\begin{itemize}
\tightlist
  \item behavioral function
    $\begin{cases}
      \text{by ad hoc} \\
      \text{by optimization}
    \end{cases}$
  \item definitive function
  \item equilibrium conditions
\end{itemize}
\end{note}
\end{blueblock}

\begin{blueblock}
\begin{note}[General way of advanced macroeconomics]
By method, there are optimization, equilibrium-conducting and comparative analysis. And all of three can do both static analysis and dynamic analysis.
\end{note}
\end{blueblock}


\begin{blueblock}
\begin{note}
Inada conditions making sure that profit maximization is feasible
\end{note}
\end{blueblock}




\newpage
\section{Keynesian Model}
In a brief way, a Keynesian model can be seen as two parts:

\begin{align}
&
\text{AS: }
\begin{cases}
Y=F(K,N) \\
\frac{W}{P}=F_{N}(K,N) \\
N_{S}=N_{s}\left(\frac{W}{P}\right)\text{,where in most cases neglected from the equations}
\end{cases}
\\&\text{AD: }
\begin{cases}
Y=C+I+G+\delta K \\
C=C(Y,T-\pi) \\
I=I(q-1), I'>0, q\equiv \frac{F_{k}-(r+\delta-\pi)}{r-\pi} \\
\frac{M}{P}=m(Y,r),m_{1}>0,m_{2}<0
\end{cases}
\end{align}


The Keynesian model consists of six equations, instead of seven in
classical model. Because labor market is not clearing. So the wage is exogenous.




\newpage
